
\section{\trUnramified{}态射}

设 $\pi\colon X\to Y$ 是概形之间的一个\trMorphism。拟凝聚层 $\Omega_\pi=\Omega_{X/Y}$ 的\trSupport{}称为\trRamificationLocus，而这个支撑的像 $\pi(\Supp\Omega_{X/Y})$ 称为\trBranchLocus。若 $\Omega_\pi=0$，我们称 $\pi$ 是\textbf{形式\trUnramified}（formally unramified）的；而若 $\pi$ 进而还是局部有限型的，我们便称 $\pi$ 是\trUnramified{}的。（警示：关于「\trUnramified」的定义尚缺乏共识；「局部有限型」有时会被替换为「局部有限表现」，后者是 [Gr-EGA] 中最初使用的定义。）

\textbf{21.7.A.（容易的练习：\trUnramified{}态射的例子）}\quad
（a）证明：\trLocallyClosedEmbedding{}都\trUnramified。
（b）证明：若 $S$ 是环 $B$ 的一个\trMultiplicativeSubset，则 $\Spec S^{-1}B\to\Spec B$ 形式\trUnramified。（于是，若 $\eta$ 是整概形 $Y$ 的\trGenericPoint，则 $\Spec\OO_{Y,\eta}\to Y$ 形式\trUnramified。）
（c）证明：有限可分\trFieldExtension（或者说得更准确些，其所对应的概形映射）都\trUnramified。

\textbf{21.7.B.（练习：概念练手）}\quad
（a）证明：\mbox{练习~10.7.E} 中那个\trNode{}的\trNormalization{}是\trUnramified{}的（见图~\ref{fig:normalization-nodal-cubic}）。
（b）证明：\mbox{练习~10.7.F} 中那个\trCusp{}的\trNormalization{}并非\trUnramified{}的（见图~\ref{fig:cusp-normalization}）。

\textbf{21.7.C.（容易的练习：\trUnramified{}态射是一类「合理」的态射）}\quad 证明：\trUnramified{}态射构成的类在 \S8.1 的意义下是「合理」的。

\textbf{21.7.D.（练习：用纤维刻画\trUnramified{}态射）}\quad 设 $\pi\colon X\to Y$ 局部有限型。
（a）证明：$\pi$ \trUnramified，当且仅当对每个 $\mathrm{q}\in Y$，$\pi^{-1}(\mathrm{q})$ 都是形如 $\Spec K$ 的概形的（概形论）不交并，其中 $K$ 是 $\kappa(\mathrm{q})$ 的一个有限可分扩张。
（b）证明：$\pi$ \trUnramified，当且仅当对每个\trGeometricPoint{} $\overline{\mathrm{q}}$，$\pi^{-1}(\overline{\mathrm{q}}):=\overline{\mathrm{q}}\times_Y X$ 都是若干份 $\overline{\mathrm{q}}$ 拷贝的（概形论）不交并。

\textbf{21.7.E.（练习：\trUnramified{}性在「非常几何」场景中的含义）}\quad 证明：代数闭域上有限型概形之间的\trMorphism{} $\pi\colon X\to Y$ \trUnramified，当且仅当 $\pi$ 在诸闭点处的切向量上是单射。（\mbox{定理~19.1.1} 正后方的那条注记曾提及此事。）

\textbf{21.7.F.（练习：供\mbox{练习~21.7.G(b)} 使用）}\quad
（a）设 $\pi\colon X\to Y$ 是局部\trNoetherian{}概形之间的一个局部有限型\trMorphism。证明：$\pi$ \trUnramified，当且仅当 $\delta_\pi\colon X\to X\times_Y X$ 是一个\trOpenEmbedding。（提示：证明如下结论。若 $\phi\colon X\to Z$ 是\trNoetherian{}概形之间的一个\trClosedEmbedding，而 $\phi$ 的\trIdealSheaf{} $\mathscr{S}$ 满足 $\mathscr{S}=\mathscr{S}^2$，则 $\phi$ 也是一个\trOpenEmbedding。为此，证明：若 $(A,\mathfrak{m})$ 是一个\trNoetherian\trLocalRing，$I$ 是 $A$ 的一个\trProperIdeal{}且满足 $I=I^2$，则 $I=0$。而为此，用\trNakayama{}引理版本~2（\mbox{引理~8.2.9}）。另外还要用到 $\Supp\mathscr{S}$ 是闭的这一事实（用\mbox{练习~14.3.C}），从而它的\trComplement{}是开集。）
\hard\ （b）把你（a）中的证明加以改造，以去掉局部\trNoetherian{}假设。（提示：证明若 $\pi$ 局部有限型，则 $\delta_\pi$ 局部有限表现。）

\textbf{21.7.G.（容易的练习）}\quad 设 $\pi\colon X\to Y$ 与 $\rho\colon Y\to Z$ 都局部有限型。令 $\tau=\rho\circ\pi$：
\begin{tikzcd}
X\ar[rr,"\pi"]\ar[dr,"\tau"'] && Y\ar[dl,"\rho"]\\
& Z.
\end{tikzcd}
（a）证明：若 $\tau$ \trUnramified，则 $\pi$ 也\trUnramified。
（b）证明：若 $\rho$ \trUnramified，而 $\tau$ 是相对维数 $n$ 的光滑\trMorphism（例如 $n=0$ 时即\trEtale），则 $\pi$ 是相对维数 $n$ 的光滑\trMorphism。
（这与你的几何直觉相合吗？）

\textbf{21.7.H.（练习）}\quad 设 $\pi\colon X\to Y$ 局部有限型。证明：$X$ 中 $\pi$ 在其上\trUnramified{}的那个轨迹是开的。

\textbf{21.7.I.（不重要的练习：致数论学家）}\quad 设 $\phi\colon(A,\mathfrak{m})\to(B,\mathfrak{n})$ 是局部环之间的一个局部同态。在代数数论中，这样的\trRingMorphism{}称为\emph{\trUnramified{}的}，若 $B/\phi(\mathfrak{m})B$ 是 $A/\mathfrak{m}$ 的一个有限可分扩张。证明：若 $\phi$ 有限型，这与我们的定义相一致。

\textbf{21.7.1.}\ \emph{算术旁注：差积与\trDiscriminant。}\quad 若 $B$ 是某个\trNumberField{}的\trRingOfIntegers{}（\S10.7.1），则 $B$ 的\textbf{\trDifferentIdeal}（different ideal）就是 $\Omega_{B/\mathbb{Z}}$ 的零化子。它度量 $\Spec B\to\Spec\mathbb{Z}$ 有多「（非）分歧」，并且是\trRamificationLocus{}的一个环论版本。\textbf{\trDiscriminantIdeal}（discriminant ideal）可以被诠释为 $\mathbb{Z}$ 中的一个理想，它对应着 $\Spec\mathbb{Z}$ 上的那条有效除子\footnote{原书此处作「corresponding to effective divisor on Spec $\mathbb{Z}$」，疑脱定冠词 the，当为「corresponding to the effective divisor on Spec $\mathbb{Z}$」，兹照录原书。——译者注}，而这条除子正是「对应于差积的那条除子」的「前推」（此处不给定义，但如你所预期的那样定义）。\trDiscriminantIdeal{}是\trBranchLocus{}的一个环论版本。若 $B/A$ 是数域的整数环之间的一个扩张，则（$B$ 的）\textbf{相对\trDifferentIdeal}（relative different ideal）与（$A$ 的）\textbf{相对\trDiscriminantIdeal}（relative discriminant ideal）可以类似地定义。（我们用不到这些想法。）
